\section{RQ3: Replacement Quality}
\label{sec:rq3}
\textit{How effectively can figurative expressions be replaced with literal, meaning-preserving alternatives?}

\subsection{Experimental Setup}

Replacement experiments use the \textit{detect-then-replace} task variant: the model detects figurative spans and in the same call produces a literal paraphrase of the full sentence. All experiments use the open-weights cluster (\texttt{meta-llama/Meta-Llama-3.1-8B-Instruct}, FP16; and \texttt{casperhansen/llama-3.3-70b-instruct-awq} as a scale comparator), temperature~0. Two prompt versions are evaluated:
\begin{itemize}
    \item \textbf{Prompt v1}: a minimal stub instructing the model to identify figurative expressions and provide a literal paraphrase (three lines of instruction).
    \item \textbf{Prompt v2}: a chain-of-thought prompt instructing the model to (1) identify each expression, (2) reason about its meaning in context, and (3) rewrite the full sentence with all figurative parts replaced, keeping non-figurative content unchanged. Includes one illustrative example.
\end{itemize}

\subsection{Idiom Replacement (SemEval 2022)}

SemEval 2022 Task B provides gold literal paraphrases for a subset of Task A sentences, enabling automatic evaluation of replacement quality using BLEU and BERTScore. Table~\ref{tab:rq3-semeval} reports detect-then-replace results across two prompts (v1 stub, v2 chain-of-thought) and two architectures (single-call structured output, 3-step pipeline) on both Llama-3.1-8B-Instruct and Llama-3.3-70B-Instruct-AWQ, on the same deterministic 200-example slice ($T=0$).

\begin{table}[h!]
\centering
\begin{tabular}{llllllr}
\hline
\textbf{Architecture} & \textbf{Prompt} & \textbf{Model} & \textbf{F1-sent} & \textbf{F1-span} & \textbf{BLEU} & \textbf{BERTScore F1} \\
\hline
Single-call CoT & v1 (stub)              & 8B  & 0.480 & 0.163 & 0.397 & 0.914 \\
Single-call CoT & v2 (chain-of-thought)  & 8B  & 0.492 & 0.162 & 0.400 & 0.915 \\
Single-call CoT & v2                     & 70B & 0.383 & 0.161 & 0.423 & 0.915 \\
\hline
3-step pipeline & ---                    & 8B  & 0.255 & 0.141 & 0.430 & 0.916 \\
3-step pipeline & ---                    & 70B & 0.435 & 0.225 & 0.414 & 0.915 \\
\hline
\end{tabular}
\caption{Idiom detect-then-replace on SemEval 2022 Task A/B, deterministic 200-example slice, $T=0$. Single-call CoT runs use the structured-output schema requiring detection + replacement in one inference; 3-step pipeline runs separate detection, explanation, and transformation into three sequential LLM calls. F1-sent and F1-span measure detection quality on the same examples. Runs: 8B v1 \texttt{b2b174ea}, 8B v2 \texttt{d6bc60b0}, 70B v2 \texttt{92941662}, 8B pipeline \texttt{3551e201}, 70B pipeline \texttt{56f19b8c}.}
\label{tab:rq3-semeval}
\end{table}

\paragraph{The chain-of-thought benefit is model-dependent and does not transfer to Llama-3.1-8B.} On 8B, the v1 and v2 prompts produce statistically indistinguishable outputs: BLEU 0.397 vs.\ 0.400 (+0.3pp), BERTScore F1 0.914 vs.\ 0.915 (+0.1pp), sentence F1 0.480 vs.\ 0.492 (+1.2pp). This contradicts a development-time A/B run on \texttt{google/gemini-3-flash-preview} which showed +11pp BLEU and +1.6pp BERTScore F1 in v2's favour for the same prompt pair (engineering log, Step 20, 100-example v1 vs.\ 20-example v2). The most parsimonious explanation is that v2's explicit identify--reason--rewrite scaffolding is benefiting closed-source large models more than the open-weights 8B Instruct, possibly because the canonical Llama Instruct chat template already implements similar internal scaffolding by default. This finding is revisited in Chapter~\ref{ch:discussion} as evidence for the \emph{model-dependence} limitation of zero-shot prompting strategies.

\paragraph{The pipeline architecture's BLEU advantage shrinks to zero at scale.} On 8B the 3-step pipeline beats single-call CoT v2 by +3pp BLEU (0.430 vs.\ 0.400). On 70B the same comparison flips slightly to a $-1$pp difference (0.414 vs.\ 0.423) --- within sample noise. BERTScore F1 is tied at $\sim$0.915 in all four cells. The pipeline's small-scale BLEU advantage is therefore concentrated at 8B and partially attributable to a known artefact: when the pipeline's detection step finds nothing, the transform step short-circuits and returns the input sentence unchanged; the input shares roughly 85\% of its tokens with the gold paraphrase, so those examples score artificially high on BLEU. This artefact contributes more to the 8B aggregate (which has lower detection recall, so more fall-throughs) than to the 70B aggregate. The cleaner read is that scale lifts detection F1 substantially while leaving replacement quality flat on automatic metrics --- a dissociation we develop in the analysis below.

\paragraph{Single-call CoT detection F1 inverts at scale, on both phenomena.} 8B SemEval idiom: detection-only sentence F1 is 0.255 (run \texttt{cfe4f798}); the same model under dtr v2 jumps to 0.492 ($+22$pp from its detection-only baseline). 70B-AWQ SemEval idiom: detection-only sentence F1 is 0.435 (run \texttt{be7dd99f}); the same model under dtr v2 \emph{drops} to 0.383 ($-5$pp from its detection-only baseline). The same prompt that helps 8B by $+22$pp hurts 70B by $-5$pp on the same dataset. Mechanism is consistent with the VU metaphor inversion (Section~\ref{sec:rq1}, $+12$pp at 8B / $-28$pp at 70B): the 70B's token precision climbs when forced to commit alongside a paraphrase, but its recall drops as it over-commits to ``no idiom'' on borderline examples. The schema scaffold helps 8B and distorts 70B.

BLEU values are low overall (0.397--0.430), as expected: valid literal paraphrases of idiomatic expressions often share no n-grams with the gold reference. The idiom \textit{smoking gun}, for example, may be correctly replaced with \textit{decisive evidence} or \textit{conclusive proof}, neither of which overlaps with the reference. BERTScore (0.914--0.916) is the more informative automatic signal here, capturing semantic proximity at the embedding level; the human-evaluation survey (Section~\ref{subsec:rq3-human-eval}) provides the definitive judgement on grammaticality, meaning preservation, and simplicity.

\subsection{Metaphor Replacement (VU Amsterdam)}

The VUAMC does not include gold literal paraphrases, making BLEU and BERTScore inapplicable. Replacement quality for metaphors is therefore assessed through (a) detection metrics on detect-then-replace runs (Table~\ref{tab:rq3-vu}) and (b) qualitative inspection. Metaphor replacement is not in scope for the human-evaluation survey: as discussed in Section~\ref{sec:eval-methodology}, the rubric triad (grammaticality, meaning preservation, simplicity) is calibrated for idiom replacement, where the figurative-to-literal mapping is well-defined; conventional metaphors (\textit{life is a journey}) often do not require replacement at all, and including them would dilute the rated pool.

\begin{table}[h!]
\centering
\begin{tabular}{lllll}
\hline
\textbf{Task} & \textbf{N} & \textbf{F1-sent} & \textbf{F1-tok} & \textbf{F1-span} \\
\hline
Detection only        & 200 & 0.430 & 0.180 & 0.052 \\
Detect-then-replace v1 & 197 & 0.553 & 0.230 & 0.067 \\
Detect-then-replace v2 & 196 & 0.551 & 0.229 & 0.065 \\
\hline
\end{tabular}
\caption{Metaphor detection metrics for detection-only vs.\ detect-then-replace runs on the deterministic 200-example VUAMC slice (Llama-3.1-8B-Instruct, $T=0$). The detect-then-replace task lifts sentence F1 by 12.1--12.3 percentage points over detection-only; v1 vs.\ v2 prompts produce indistinguishable results (within 0.2pp on every metric).}
\label{tab:rq3-vu}
\end{table}

\subsection{Analysis}

\paragraph{BLEU is structurally misaligned with replacement evaluation.} For tasks with multiple valid outputs, BLEU's n-gram overlap requirement penalises correct paraphrases that diverge lexically from the reference. BERTScore, which compares contextual embeddings, partially addresses this; but neither correlates reliably with human judgement on simplification-style tasks~\cite{Alva-Manchego2021,Scialom2021}, motivating the human-evaluation protocol below.

\paragraph{Idioms outperform metaphors at every replacement-relevant metric (H3).} Direct sentence-F1 comparison: 0.480--0.492 (idiom dtr) vs.\ 0.551--0.553 (metaphor dtr) at first glance favours metaphors, but this conflates dataset granularity (token-level VU annotation vs.\ span-level SemEval) with phenomenon difficulty. The cleaner indicator is that on idioms the system produces a meaning-preserving literal paraphrase (BLEU 0.40, BERTScore 0.91 against gold) whereas the VUAMC has no gold paraphrase against which to score the metaphor rewrite at all --- the absence of a measurable replacement metric for metaphors itself is consistent with H3's underlying claim that metaphor-to-literal mappings are less constrained. The human-evaluation survey covers idiom replacement only and will provide the definitive H3 evidence on the rubric dimensions.

\paragraph{The detect-then-replace prompt produces stronger detection than detection-alone on Llama-3.1-8B-Instruct.} On both datasets, sentence F1 in detect-then-replace exceeds detection-only by 12 percentage points (VU: 0.430 $\to$ 0.553; SemEval: 0.255 $\to$ 0.480). This is the opposite of the development-time pattern observed under hosted closed-source models, where adding a replacement requirement degraded span localisation. On the open-weights system, the structured-output schema in detect-then-replace appears to disambiguate the task for the model: forcing it to commit to a specific paraphrase pulls along sharper span identification. This is itself early evidence for RQ2 (decomposition $>$ monolithic): even the trivial decomposition of \emph{adding the replacement step} changes detection behaviour for the better at this scale.

\paragraph{Scale (8B $\to$ 70B) lifts detection but not replacement, confirmed at 200 examples.} The 30-source survey-pool comparison (single-call agentic v2: 8B sentence F1 0.609 $\to$ 70B 0.767, $+15.8$pp; BLEU $0.061 \to 0.058$; BERTScore $0.851 \to 0.858$) and the 200-example deterministic slice (Table~\ref{tab:rq3-semeval}: pipeline architecture sentence F1 $0.255 \to 0.435$, $+18$pp; BLEU $0.430 \to 0.414$; BERTScore $0.916 \to 0.915$) both show the same dissociation: detection sentence F1 lifts $+16$ to $+18$ percentage points with scale, but BLEU and BERTScore are tied or slightly lower at 70B. The within-family scale comparator therefore confirms a clean RQ1 finding: at this size class, scale lifts the model's ability to \emph{find} the figurative expression but not its ability to \emph{rewrite} it. Whether the survey's rubric ratings (grammaticality, meaning preservation, simplicity) confirm or contradict this dissociation will be reported in Section~\ref{subsec:rq3-human-eval-results}.

\subsection{Human Evaluation}
\label{subsec:rq3-human-eval}

The replacement task is evaluated by Direct Assessment~\cite{Graham2013,Graham2017}; the protocol, rubric, standardisation, and three departures from the original DA specification are described in Chapter~\ref{ch:methodology}. This section gives the survey \textit{instantiation} --- conditions compared, source-sentence pool, item counts, and respondent recruitment. The full instrument (item screenshots, consent form, bad-reference construction recipe) is reproduced in Appendix~\ref{appendix:survey-design}.

\paragraph{Conditions.} The survey rates outputs from three open-weights conditions, all on the same set of source sentences:
\begin{itemize}
    \item \textbf{Llama-3.1-8B + agentic detect-then-replace} --- the deliverable system;
    \item \textbf{Llama-3.1-8B + monolithic single-prompt} --- the RQ2 ablation, isolating the contribution of agentic decomposition;
    \item \textbf{Llama-3.3-70B-Instruct-AWQ + agentic detect-then-replace} --- the within-family scale comparator for RQ1.
\end{itemize}
No hosted-LLM condition is included in the human evaluation: the survey is reserved for the open-weights deliverable system and its controlled variants.

\paragraph{Pool and instrument.} Sources are drawn from a candidate pool of 50 SemEval-idiom items selected from a Llama-3.1-8B detect-then-replace run, filtered for sentence length ($<30$ words), expression length ($\leq 4$ words), and presence of a system-produced replacement; 30 sources are sampled for the survey. Each source is rated under all three conditions, yielding $30 \times 3 = 90$ items. With $\geq 3$ ratings per item and 16 rated items + 4 quality-control items per respondent, the protocol requires approximately 17 native-English-speaker respondents.

\paragraph{Rubric.} Three independent 0--100 continuous sliders per item: grammaticality, meaning preservation, simplicity (Section~\ref{sec:eval-methodology}). Annotators are blind to the system-of-origin; system identifiers live in analysis-side metadata only.

\paragraph{Quality control.} Two bad-reference items and two repeat-pair items per respondent, embedded among the 16 rated items (matching the 20\% QC ratio of \cite{Graham2013}); fixed-threshold filters are used in place of \cite{Graham2017}'s per-worker significance test, for the small-sample reasons documented in Section~\ref{sec:qc-filter-design}.

\paragraph{Aggregation.} Per-annotator z-score standardisation before per-item aggregation, with mean $\pm$ 95\% confidence interval reported for each rubric dimension and a composite summary across the three. Agreement diagnostics include intra-annotator $\kappa$ from repeat pairs and inter-annotator $\kappa$ across items.

% TODO: Add qualitative examples of correct, partially correct, and incorrect replacements once the open-weights runs land.

\newpage
